\begin{statementbox}
\textbf{几何定义。} 设 $F$ 不在准线 $l$ 上，$P$ 在非退化圆锥曲线上，$H$ 是 $P$ 到 $l$ 的垂足，$PH=d(P,l)>0$。
对\textbf{相互对应}的焦点、准线，恒定比值
\[
 \boxed{e=\frac{PF}{PH}}
\]
称为离心率。椭圆、双曲线的代数公式为
\[
 \boxed{e=\frac ca},\qquad
 \begin{cases}
 0<e<1,&\text{椭圆，}\ a>b>0,\ c^2=a^2-b^2,\\
 e>1,&\text{双曲线，}\ a,b>0,\ c^2=a^2+b^2.
 \end{cases}
\]
$a$ 为长半轴或实半轴，$b$ 为短半轴或虚半轴，$c$ 为半焦距。
抛物线由 $PF=PH$ 直接得到 $\boxed{e=1}$，不能套用 $c/a$。
\[
 e_{\text{椭}}=\sqrt{1-b^2/a^2},\qquad
 e_{\text{双}}=\sqrt{1+b^2/a^2}.
\]
中心在原点、焦点在 $x$ 轴时，$(c,0)$ 与 $x=a^2/c=a/e$ 对应，负号一侧同理；纵向时交换 $x,y$。
圆可另约定 $e=0$，但不属于这里有限焦点、有限准线的距离比模型。


\medskip
\textbf{考场抓手。} 准线距离是点到直线的\textbf{垂直距离}，不是点到准线上任一点的距离；
$e$ 是\textbf{同一焦点配同侧对应准线}得到的恒定比例。
求数值时首选 $e=c/a$；若要从头说明为什么成立，可以从第一定义推出
$PF_R=|a-(c/a)x|$，再读出准线 $x=a^2/c$。
\textbf{最少记忆：} 不必死记准线位置，也可以用左右两个顶点的焦准比消元推出 $e=c/a$（证明详见后文）。
\end{statementbox}
